\chapter{Mobilna aplikacija}

Mobilna aplikacija je jedinstvena Flutter kodna baza koja se, u zavisnosti od uloge prijavljenog korisnika (poglavlje 7.3), grana na dva različita korisnička toka: vozača i dispečera. Grananje se vrši na jednom mestu (\texttt{entry\_router}), odmah nakon prijave, na osnovu polja \texttt{role} vraćenog sa backend-a.

\section{Tok korisnika — vozač}

Nakon prijave (\texttt{login\_screen}, \texttt{register\_screen} sa izborom uloge vozač/dispečer preko \texttt{SegmentedButton}-a, ili \texttt{forgot\_password\_screen} za obnovu lozinke), samostalni vozač ili vozač koji radi za dispečera dolazi na svoj početni ekran. Ako je trenutno pod dispečerom i čeka na ponuđenu turu, prikazuje mu se \texttt{offered\_trips\_screen} sa ponudama koje treba prihvatiti ili odbiti; u suprotnom, vozač sam pokreće planiranje preko \texttt{route\_request\_screen}.

\textbf{Planiranje rute.} \texttt{route\_request\_screen} omogućava izbor polazišta (trenutna GPS pozicija ili pretraga adrese preko \texttt{address\_search\_field}, koja poziva backend \texttt{GET /geocode} proxy ka Nominatim-u) i odredišta, izbor vozila (preko \texttt{vehicle\_list\_screen}/\texttt{vehicle\_profile\_screen}, gde vozač unosi ili menja visinu, širinu, dužinu, masu, osovinsko opterećenje i prevoz opasnog tereta), i prikaz vraćenih alternativnih ruta na mapi, sortiranih po riziku (poglavlje 6.2), sa mogućim objašnjenjem odstupanja rute (poglavlje 6.3) prikazanim kao tekstualna napomena.

\textbf{Pokretanje i praćenje ture.} Pre pokretanja ture, \texttt{start\_proximity\_status} widget provera da je vozač stvarno u blizini (unutar 500m) prijavljenog polazišta preko GPS-a — namerna provera koja sprečava pokretanje ture "iz kancelarije" i garantuje da je prva pozicija koju backend primi zaista stvaran GPS fix, ne proizvoljna vrednost (relevantno za odluku da se ukloni simulacija pozicije, poglavlje 7.5). Tokom vožnje, \texttt{active\_trip\_screen} prikazuje trenutnu poziciju, procenjeno vreme dolaska i alert za predlog pauze (poglavlje 7.4) primljene preko WebSocket veze (\texttt{trip\_socket.dart}, sa automatskim ponovnim povezivanjem po rastućem intervalu — 1, 2, 5, potom 10 sekundi — u slučaju privremenog prekida konekcije), a \texttt{active\_trip\_banner} ostaje vidljiv i na drugim ekranima aplikacije dok je tura aktivna, kao podsetnik i prečica nazad na \texttt{active\_trip\_screen}. Ako vozač napusti planiranu rutu za više od 300m, backend automatski preračunava novu rutu od trenutne pozicije do istog odredišta (\texttt{POST /trips/\{id\}/reroute}).

\textbf{Dnevnik i istorija.} \texttt{trip\_log\_screen} prikazuje dnevnik događaja jedne ture (\texttt{trip\_events} — polazak, predlog pauze, dolazak) kao vremensku liniju, dok \texttt{trip\_list\_screen} organizuje sve vozačeve ture u tri kategorije (u toku, predstojeće, završene). \texttt{truck\_status\_screen} prikazuje nivo goriva, preostali kilometre do sledećeg servisa i ukupne sate vožnje vozila; \texttt{preferences\_screen} omogućava podešavanje stila vožnje (klizači za prioritet potrošnje goriva, osetljivosti tereta, korišćenja auto-puta i brzine, opisani u poglavlju 6.2) i omiljene benzinske stanice; \texttt{chat\_list\_screen}/\texttt{chat\_thread\_screen} omogućavaju dopisivanje sa dispečerom u realnom vremenu (poglavlje 7.5).

\section{Tok korisnika — dispečer}

Dispečer upravlja flotom vozača i vozila kroz odvojen skup ekrana. \texttt{dispatcher\_home\_screen} je početna tačka; \texttt{dispatcher\_available\_drivers\_screen} prikazuje vozače koji trenutno nisu vezani za nijednog dispečera i omogućava slanje zahteva za povezivanje (poglavlje 7.3), dok \texttt{dispatcher\_requests\_screen} prikazuje status poslatih zahteva. \texttt{dispatcher\_create\_trip\_screen} omogućava kreiranje i dodelu ture jednom od svojih vozača, birajući između vozačevog sopstvenog vozila ili vozila u vlasništvu dispečera (poglavlje 7.3).

Najkompleksniji ekran dispečerskog toka je \texttt{dispatcher\_live\_map\_screen}: prikazuje jednu mapu sa \textbf{N paralelnih WebSocket konekcija}, po jednu za svaku aktivnu turu svojih vozača, i ažurira poziciju svakog vozila na mapi čim gateway (poglavlje 7.5) emituje novu poziciju za odgovarajuću turu. \texttt{dispatcher\_trip\_list\_screen} prikazuje istoriju svih tura koje je dispečer dodelio.

\section{Arhitektura mobilne aplikacije}

Aplikacija je organizovana u standardnu Flutter strukturu: \texttt{screens/} (ekrani navedeni u 8.1/8.2), \texttt{models/} (podatkovne klase koje odgovaraju backend JSON odgovorima — \texttt{driver}, \texttt{vehicle\_profile}, \texttt{trip}, \texttt{trip\_event}, \texttt{route\_result}/\texttt{route\_candidate}, \texttt{chat\_message}, \texttt{dispatcher\_request}, \texttt{driver\_preferences}, \texttt{favorite\_stop}, \texttt{geocode\_result}, \texttt{position\_update}, \texttt{rest\_stop}, \texttt{account\_status}, \texttt{auth\_result}), \texttt{services/} i \texttt{widgets/} (deljene, ponovo iskoristive komponente).

Sloj servisa uključuje: \texttt{api\_client.dart} (centralizovan REST klijent koji, pored slanja zahteva, globalno presreće HTTP 401 odgovor i automatski prebacuje korisnika na ekran za prijavu, umesto da svaki ekran posebno rukuje isticanjem sesije), \texttt{auth\_storage.dart} (lokalno, trajno čuvanje JWT tokena na uređaju), \texttt{chat\_socket.dart} i \texttt{trip\_socket.dart} (WebSocket klijenti sa automatskim ponovnim povezivanjem), \texttt{location\_service.dart} (očitavanje GPS pozicije preko \texttt{geolocator} biblioteke), \texttt{google\_auth.dart} (Google prijava na klijentskoj strani), \texttt{polyline.dart} (dekodiranje enkodovane geometrije rute za prikaz na mapi) i \texttt{route\_observer.dart} (Flutter \texttt{RouteAware} mehanizam koji osvežava podatke na ekranu kada se korisnik na njega vrati nakon zatvaranja nekog drugog ekrana — npr. povratak na listu vozila nakon uređivanja profila vozila).

Među deljenim widgetima ističe se \texttt{radial\_fab\_menu.dart} — radijalni, povlačivi (\emph{draggable}) meni pokrenut iz plutajućeg dugmeta (\emph{floating action button}), korišćen kao glavna navigacija umesto klasične donje trake sa karticama (\emph{bottom navigation bar}), kao i \texttt{email\_verification\_banner} (podsetnik da nalog nije verifikovan) i \texttt{start\_proximity\_status} (opisan u 8.1).

\clearpage
